% R06, 2026-09-20. Draft fragment; not merged into the original manuscript.
% Requires the family and the centered coordinates to be defined earlier.
% No priority claim is made; bibliography is in references.bib.

\paragraph{Relation to earlier work.}
The deformation considered here contains the classical de Bruijn--Newman
family: in the normalization of Rodgers and Tao \cite{RodgersTao2020},
\(F(t;\lambda,0,0)=H_\lambda(2t)\) and
\(F(t;0,0,0)=\xi(1/2+it)/8\).
For fixed real \(\mu,\nu\), the positive symmetric measure
\(d\rho_{\mu,\nu}(u)=\Phi(|u|)e^{\mu u^4+\nu u^6}\,du\)
places this family within the measure-based setting considered by
Newman and Wu \cite{NewmanWu2020}.
Our results concern the local discriminant of this particular
multi-parameter deformation, rather than the condition that all complex
zeros of a transform are real. In particular, they give no new bound
for the classical de Bruijn--Newman constant.

The two folds and the one-versus-three local zero configuration are
standard cusp geometry. Computer-assisted verification of cusp
bifurcations in ODEs is developed by Lessard and Pugliese
\cite{LessardPugliese2025}. Validated multi-parameter continuation and
gluing of local charts are established tools \cite{GameiroLessardPugliese2016};
continuous branches of fold and pitchfork points have also been
rigorously computed \cite{LessardSanderWanner2017}.
Accordingly, we do not claim a new general cusp classification or a new
general continuation method. The PDE work of Blanco and Lessard
\cite{BlancoLessard2026}, currently a preprint, further illustrates the
distinction between a local multiplicity picture and dynamical
stability, which in that setting requires additional spectral information.

Heat-equation identities also connect our problem to existing zero
theory. At fixed \(\mu,\nu\), setting \(\lambda=-4\tau\) gives
\(\partial_\tau F=\partial_t^2F\).
The isolation of multiple zeros of nontrivial entire caloric functions
is proved by Papanicolaou, Kallitsi, and Smyrlis
\cite[Theorem~4.1]{PapanicolaouKallitsiSmyrlis2021}.
This does not prevent a cusp curve when the additional deformation
parameters vary. It also does not determine the specific connected
arc, its transverse zero regions, or their metric variation studied here.

The specific contribution is a computer-assisted description of an
analytic cusp arc joining the stated quartic and sextic sections,
together with uniform real-zero counts and strictly nested centered
three-zero regions. The latter comparison is made in the original,
unrescaled monomial control coefficients, for every
\(\nu\in[-29,0]\) and \(0<\ell\le10^{-6}\), with explicit bounds on
\(-\partial_\nu W\).
The opening coefficient is coordinate dependent. Moreover, once its
strict derivative sign is known on a compact analytic arc, qualitative
nesting in some sufficiently small common neighborhood follows by
analyticity. The finite-window theorem supplies a specified window
and quantitative remainder control; it is not a separate universal
cusp law. All zero counts refer to the stated local \(t\)-interval,
and nesting refers to coordinates centered at each exact cusp.
